\RequirePackage[nomain,symbols,acronym,xindy={language=german-din},toc,nostyles]{glossaries}
\RequirePackage{tabularx}
\RequirePackage{ltablex}
\RequirePackage{booktabs}
\RequirePackage{environ}
\RequirePackage{expl3}
\RequirePackage{xparse}
\glsdisablehyper
\keepXColumns              

%%SYMBOLVERZEICHNIS%%

\newcommand{\fixgap}{\noalign{\vskip-\bigskipamount}\noalign{\vskip-\medskipamount}}


%%FORMELZEICHEN%%
\newglossaryentry{Formelzeichen}{%
type=symbols,%
name={\bfseries{Skalare und Vektoren}},%
description={\nopostdesc},%
sort = {a}
}

%%MATRIZEN%%
\newglossaryentry{matrices}{%
type=symbols,%
name={\bfseries{Matrizen}},%
description={\nopostdesc},%
sort = {b}
}


%%INDICES%%
\newglossaryentry{Indices}{%
type=symbols,%
name={\bfseries{Indizierungen}},%
description={\nopostdesc},%
sort = {c}
}


\ExplSyntaxOn

% Central variable for sort key manipulation
\tl_new:N \l_autosort_tl

% --- CENTRAL GREEK DICTIONARY ---
% Replaces LaTeX macros with numeric prefixes to enforce DIN-sorting
\cs_new_protected:Nn \apply_greek_sorting:N
  {
    \tl_replace_all:Nnn #1 { \alpha } { 01alpha }
    \tl_replace_all:Nnn #1 { \beta  } { 02beta }
    \tl_replace_all:Nnn #1 { \gamma } { 03gamma }
    \tl_replace_all:Nnn #1 { \Gamma } { 03Gamma }
    \tl_replace_all:Nnn #1 { \delta } { 04delta }
    \tl_replace_all:Nnn #1 { \Delta } { 04Delta }
    \tl_replace_all:Nnn #1 { \epsilon } { 05epsilon }
    \tl_replace_all:Nnn #1 { \varepsilon } { 05varepsilon }
    \tl_replace_all:Nnn #1 { \zeta } { 06zeta }
    \tl_replace_all:Nnn #1 { \eta } { 07eta }
    \tl_replace_all:Nnn #1 { \theta } { 08theta }
    \tl_replace_all:Nnn #1 { \Theta } { 08Theta }
    \tl_replace_all:Nnn #1 { \vartheta } { 08vartheta }
    \tl_replace_all:Nnn #1 { \iota } { 09iota }
    \tl_replace_all:Nnn #1 { \kappa } { 10kappa }
    \tl_replace_all:Nnn #1 { \lambda } { 11lambda }
    \tl_replace_all:Nnn #1 { \Lambda } { 11Lambda }
    \tl_replace_all:Nnn #1 { \mu } { 12mu }
    \tl_replace_all:Nnn #1 { \nu } { 13nu }
    \tl_replace_all:Nnn #1 { \xi } { 14xi }
    \tl_replace_all:Nnn #1 { \Xi } { 14Xi }
    \tl_replace_all:Nnn #1 { \pi } { 15pi }
    \tl_replace_all:Nnn #1 { \Pi } { 15Pi }
    \tl_replace_all:Nnn #1 { \rho } { 16rho }
    \tl_replace_all:Nnn #1 { \varrho } { 16varrho }
    \tl_replace_all:Nnn #1 { \sigma } { 17sigma }
    \tl_replace_all:Nnn #1 { \Sigma } { 17Sigma }
    \tl_replace_all:Nnn #1 { \tau } { 18tau }
    \tl_replace_all:Nnn #1 { \upsilon } { 19upsilon }
    \tl_replace_all:Nnn #1 { \Upsilon } { 19Upsilon }
    \tl_replace_all:Nnn #1 { \phi } { 20phi }
    \tl_replace_all:Nnn #1 { \Phi } { 20Phi }
    \tl_replace_all:Nnn #1 { \varphi } { 20varphi }
    \tl_replace_all:Nnn #1 { \chi } { 21chi }
    \tl_replace_all:Nnn #1 { \psi } { 22psi }
    \tl_replace_all:Nnn #1 { \Psi } { 22Psi }
    \tl_replace_all:Nnn #1 { \omega } { 23omega }
    \tl_replace_all:Nnn #1 { \Omega } { 23Omega }
  }

% --- INTERNAL WRAPPERS ---
% These macros handle the actual glossaries entry creation with forced expansion (V-variant)

\cs_new_protected:Nn \define_formula_entry:nnnn
  {
    \newglossaryentry{#1}{
      type=symbols, name={#1}, description={\nopostdesc},
      symbol={\ensuremath{#2}}, user1={\ensuremath{\left[\mathrm{#3}\right]}},
      sort={#4}, parent={Formelzeichen}
    }
  }
\cs_generate_variant:Nn \define_formula_entry:nnnn { nnnV }

\cs_new_protected:Nn \define_matrix_entry:nnnn
  {
    \newglossaryentry{#1}{
      type=symbols, name={#1}, description={\nopostdesc},
      symbol={\ensuremath{\mathbf{#2}}}, user1={\ensuremath{\left[\mathrm{#3}\right]}},
      sort={#4}, parent={matrices}
    }
  }
\cs_generate_variant:Nn \define_matrix_entry:nnnn { nnnV }

\cs_new_protected:Nn \define_indice_entry:nnn
  {
    \newglossaryentry{#1}{
      type=symbols, 
      name={#1}, 
      description={\nopostdesc},
      symbol={\ensuremath{()\sb{\mathrm{#2}}}}, %
      user1={\ensuremath{\left[\mathrm{-}\right]}},
      sort={#3}, 
      parent={Indices}
    }
  }
\cs_generate_variant:Nn \define_indice_entry:nnn { nnV }

% --- USER COMMANDS ---

% 1. Standard Formula Symbols
\NewDocumentCommand{\newformulasymbol}{ O{} m m m }
  {
    \tl_if_blank:nTF { #1 }
      { \tl_set:Nn \l_autosort_tl { #3 } \apply_greek_sorting:N \l_autosort_tl \tl_put_right:Nn \l_autosort_tl { #2 } }
      { \tl_set:Nn \l_autosort_tl { #1#2 } }
    \define_formula_entry:nnnV { #2 } { #3 } { #4 } \l_autosort_tl
  }

% 2. Matrices (bold symbol)
\NewDocumentCommand{\newmatrix}{ O{} m m m }
  {
    \tl_if_blank:nTF { #1 }
      { \tl_set:Nn \l_autosort_tl { #3 } \apply_greek_sorting:N \l_autosort_tl \tl_put_right:Nn \l_autosort_tl { #2 } }
      { \tl_set:Nn \l_autosort_tl { #1#2 } }
    \define_matrix_entry:nnnV { #2 } { #3 } { #4 } \l_autosort_tl
  }

% 3. Indices
\NewDocumentCommand{\newindice}{ O{} m m }
  {
    \tl_if_blank:nTF { #1 }
      { \tl_set:Nn \l_autosort_tl { #3 } \apply_greek_sorting:N \l_autosort_tl \tl_put_right:Nn \l_autosort_tl { #2 } }
      { \tl_set:Nn \l_autosort_tl { #1#2 } }
    \define_indice_entry:nnV { #2 } { #3 } \l_autosort_tl
  }

\ExplSyntaxOff


%%Verhaltensdefinition im Dokument (Mathe und Text) 

\defglsentryfmt[symbols]{%
\ifmmode%
\glssymbol{\glslabel}%
\else%
\glsgenentryfmt~\glsentrysymbol{\glslabel}% %% Alternativen Label oder Label Symbol 
\fi%
}


%Definition Stil Symbolverzeichnis als Tabelle mit Tablurax und Environ



\NewEnviron{tableenv}{%
\renewcommand{\arraystretch}{1.2} %Zeilenabstand
  \begin{tabularx}{\columnwidth}{l@{\hspace{3em}} l@{\hspace{3em}} X} % Änderungen des Tabellendesign hier möglich 
      \BODY
  \end{tabularx}
}


\newglossarystyle{symboltablr}
{%
        \newcommand*\symbollevel{-1}
  \renewenvironment{theglossary}%
    {\tableenv}%
    {\endtableenv}%
  
  \renewcommand*{\glsgroupheading}[1]{}%
  \renewcommand*{\glsgroupskip}{}%

 \newcommand*\symbolhead{%
\toprule
\bfseries Symbol & \bfseries Einheit &
\bfseries Bedeutung \\ \midrule
}%

\renewcommand*{\glossaryheader}{ \endfirsthead
\symbolhead \endhead
\midrule \endfoot
\bottomrule \endlastfoot
\gdef\symbollevel{-1}%
  }

  
\renewcommand*{\glossentry}[2]{%
\ifglshaschildren{##1}{%
\ifnum\symbollevel>0\relax%
\tabularnewline\fixgap\bottomrule\tabularnewline[\smallskipamount]%%\bottomrule%
\fi
\gdef\symbollevel{0}
\tabularnewline[%
        \arraystretch\dimexpr-\ht\strutbox-\dp\strutbox\relax%
      ]%
\multicolumn{3}{l}{\glossentryname{##1}}}
}%

\renewcommand*{\subglossentry}[3]{%
\ifnum\symbollevel=0\relax%
\tabularnewline[\medskipamount]\symbolhead%
\fi%
\gdef\symbollevel{##1}%
\glstarget{##2}{\glossentrysymbol{##2}} &
\glsentryuseri{##2} &
\glossentryname{##2} \tabularnewline
}%
}


%%ENDE DEFINITON



%%Definition Abkürzungsverzeichnis

\NewEnviron{acroenv}{%
\renewcommand{\arraystretch}{1.2} %Zeilenabstand
  \begin{tabularx}{\columnwidth}{l@{\hspace{4em}} X} % Änderungen des Tabellendesign hier möglich 
      \BODY
  \end{tabularx}
}

\newglossarystyle{acrotab}
{%
  \renewenvironment{theglossary}%
    {\acroenv}%
    {\endacroenv}%
  
  \renewcommand*{\glsgroupheading}[1]{}%
  \renewcommand*{\glsgroupskip}{}%

 \newcommand*\acrohead{%
\toprule
\bfseries Abkürzung & \bfseries Bedeutung \\ \midrule
}%

\renewcommand*{\glossaryheader}{
\acrohead \endfirsthead
\acrohead \endhead
\midrule \endfoot
\bottomrule \endlastfoot}
  
\renewcommand*{\glossentry}[2]{%
\glsentryitem{##1}% Entry number if required
\glstarget{##1}{\glossentryname{##1}} & \glsentrydesc{##1} \\
}
}
%%ENDE DEFINITION

%%Alternative Abkürzungsverzeichnis
\newglossarystyle{acrolist}{%
  \renewenvironment{theglossary}%
    {\begin{description}}%
    {\end{description}}%
  \renewcommand*{\glsgroupheading}[1]{}%
  \renewcommand*{\glsgroupskip}{}%
  \renewcommand*{\glossentry}[2]{%
    \item[\glsentryitem{##1}\glstarget{##1}{\glossentryname{##1}}] \glsentrydesc{##1}%
  }%
}

